\chapter{Batch Clearing Algorithm}
\label{app:algorithm}

The clearing routine implements \eqref{eq:demandsupply} and
\eqref{eq:clearing} by a search over the admissible price grid, followed
by the tie-break cascade and the rationing step of
\Cref{sec:rationing}.

\begin{lstlisting}[language=,caption={Batch clearing, single asset
pair.},label={lst:clearing}]
INPUT : buys B = {(pi_i, q_i)}, sells S = {(pi_j, q_j)},
        tick grid P, tie-break cascade C, rationing rule R,
        previous clearing price p_prev
OUTPUT: clearing price p*, fills F, residual E

1  candidates <- distinct limit prices in B and S, snapped to P
2  for each p in candidates:
3      D(p) <- sum of q_i over buys with pi_i >= p
4      S(p) <- sum of q_j over sells with pi_j <= p
5      V(p) <- min(D(p), S(p))
6  Vmax  <- max over p of V(p)
7  M     <- { p : V(p) = Vmax }              // maximising set
8  if Vmax = 0: return (no clearing, all orders roll over)
9  M     <- argmin over p in M of |D(p) - S(p)|        // cascade 1
10 if |M| > 1: M <- argmin over p in M of |p - p_prev|  // cascade 2
11 if |M| > 1: M <- argmin over p in M of p             // cascade 3
12 p*    <- unique element of M
13 fill all buys with pi_i > p* completely            // inframarginal
14 fill all sells with pi_j < p* completely
15 marginal <- orders with pi = p* on the long side
16 allocate Vmax - (filled inframarginal) among marginal using R
17 E     <- unmatched quantity on the long side
18 apply time-in-force to E: roll over, or cancel
19 return (p*, F, E)
\end{lstlisting}

Steps 1--7 are the line search; the unimodality of $V$ noted in
\Cref{sec:clearing} guarantees that the maximising set is an interval on
the grid. Steps 9--11 are the tie-break cascade: minimise the residual
imbalance, then minimise distance from the auction's own previous
clearing price $p_{\mathrm{prev}}$ (not an external reference — the FBA
has no oracle input, see \Cref{app:config}), then take the lowest price
as a final deterministic tie-break. Step 16 is the only place where the
rationing rule enters, which is why the rule can be swapped without
touching the price computation.

\chapter{Engine Configuration Parameters}
\label{app:config}

\begin{table}[htbp]
\centering
\caption[Engine configuration]{The engines' fixed configuration and
behaviour, as actually implemented. This thesis used one configuration
throughout --- most rows below are not runtime-adjustable in the current
codebase, which the Meaning column states plainly rather than presenting
a wider configuration surface than exists.}
\label{tab:config}
\small
\begin{tabularx}{\textwidth}{@{}p{3.2cm} L p{3.2cm}@{}}
\toprule
\textbf{Parameter} & \textbf{Meaning} & \textbf{Value}\\
\midrule
\texttt{engine}            & CDA or FBA & both, paired\\
\texttt{tau\_ns}           & Batch interval, in nanoseconds; the same
value also sets the metrics CSV's bucket width (\Cref{sec:environment})
& 1{,}000{,}000{,}000 (1\,s)\\
\texttt{boundary}          & Batch-close timing & fixed only --- no
randomised-boundary code path exists\\
\texttt{rationing}         & Allocation of the marginal quantity at
$\clr$ & fixed: price-time-within-the-batch (rank by aggressiveness,
then submission time, then order id) --- not pro rata, not a runtime
choice\\
\texttt{tie\_break}        & Cascade for selecting $\clr$ from the
volume-maximising set & fixed: (1)~maximise matched volume, (2)~minimise
residual imbalance, (3)~minimise distance from the auction's own
previous clearing price, (4)~lowest price\\
\texttt{residual\_policy}  & Disposition of unmatched quantity at $\clr$
& fixed: always rolled into the next batch --- no cancel-on-no-clear
path exists\\
\bottomrule
\end{tabularx}
\end{table}

Four rows present in an earlier draft of this table --- a seed for a
randomised boundary, a $\kappa$ parameter for a size/pro-rata hybrid
rationing rule, a self-trade-prevention policy, and a warm-up period
excluded from metrics --- have been removed rather than marked
\texttt{n/a}: none of the four has any basis in the current codebase
(no randomised boundary, no hybrid rationing rule, no self-trade
prevention, and no warm-up exclusion are implemented), and a table row
implies a real, if unused, feature that keeping them here would have
been misleading. Self-trade prevention's absence in particular is a
genuine limitation, not a controlled-for policy shared by both engines
--- it is stated as such in \Cref{sec:experiments}.

\chapter{Reproducibility}
\label{app:repro}

Every run's checkpoint file records the configuration of
\Cref{tab:config} together with the input source path, the interval
$\tau$, and per-engine running totals, so that a resumed run can be
checked against the run it continues. A per-file input hash and a code
commit hash embedded in the run's own output --- which would let a
result be tied unambiguously to the exact code and data that produced
it, beyond what the configuration and file list already establish ---
are not currently written by the checkpoint format; this is listed here
as a gap worth closing, not implied to already exist. Given identical
input and configuration, an uninterrupted run reproduces byte-identical
output (\Cref{sec:experiments}). The analysis pipeline reads only the
logged CSV output and is shared by both engines, so no metric is
computed by engine-specific code.

\ph[repository URL, dataset DOI and commit hash to be inserted before
submission]

\chapter{Metric Formulas}
\label{app:formulas}

This appendix gives the exact, implementation-level formula behind every
column of the logged output that \Cref{sec:metrics} summarises, and
behind the handful of further diagnostic columns computed alongside them
but not tabulated in \Cref{tab:metrics1,tab:metrics2,tab:metrics3}. Each
entry carries a scope tag: \textbf{U}niversal (both engines),
\textbf{C}DA-only, \textbf{F}BA-only, or \textbf{X} (defined but always
empty --- the dataset lacks what it would need). A metric marked
\textbf{X} is kept in the catalogue rather than silently dropped, so
that its absence from every result in \Cref{ch:results} is a documented
property of the data, not an omission.

\section{Bucketing and Notation}

Every column is aggregated into fixed-width time buckets:
\begin{equation}
  \mathrm{bucket}(t) = \mathrm{anchor} + \left\lfloor \frac{t -
  \mathrm{anchor}}{\tau} \right\rfloor \tau,
\end{equation}
where $\mathrm{anchor}$ is the timestamp of the first record of the run
and $\tau$ is the interval (1\,second throughout this thesis,
\Cref{app:config}). Buckets are contiguous and empty-row-inclusive: a
bucket with no qualifying observation still produces a row, so a quiet
period is a gap in the data, not a gap in the file.

A missing value (\emph{None}, an empty CSV field) and an actual zero are
kept distinct and mean different things: \emph{None} is ``no qualifying
observation this bucket,'' zero is ``there were observations and the
value is genuinely zero.'' Every metric below follows this convention
with one deliberate exception, the intra-interval price dispersion
entry of \Cref{sec:formularq22} below. Price- and notional-denominated raw
values are stored in a fixed-point integer scale,
$\mathrm{PRICE\_SCALE} = 10^{6}$; bps and ratio metrics are already
unit-free and are \emph{not} rescaled.

\section{RQ2.1: Liquidity and Transaction Cost Metrics}

\paragraph{Quoted spread (U).} CDA: for every book snapshot in the
bucket with both sides present and $a_t + b_t > 0$,
$\mathrm{spread}_{\mathrm{bps}} = (a_t - b_t)/\midp \cdot 10^{4}$,
averaged over qualifying snapshots. FBA: for every batch with both a
best unfilled buy and a best unfilled sell and $\clr > 0$,
$(\text{unfilled sell} - \text{unfilled buy})/\clr \cdot 10^{4}$,
averaged over qualifying batches. A snapshot or batch that fails its own
gate is excluded from the mean, not counted as zero.

\paragraph{Depth at best (U).} CDA: mean of (best bid size $+$ best ask
size)$/2$ across snapshots --- touch only, not the whole book. FBA: mean
of (demand at $\clr$ $+$ supply at $\clr$)$/2$ across batches.

\paragraph{Depth within $x$ bps (U), $x \in \{10, 50, 100\}$.} For every
resting order, $\mathrm{bps\_away} = |p - \mathrm{ref}| \cdot
10^{4}/\mathrm{ref}$ against the snapshot midpoint or batch clearing
price; its quantity accumulates into every threshold with
$\mathrm{bps\_away} \le x$ --- the three columns are cumulative, so a
level within 10\,bps also counts toward the 50 and 100\,bps columns.

\paragraph{Book imbalance (C).} $(V^{bid} - V^{ask})/(V^{bid}+V^{ask})$
at the top of book, averaged over snapshots with nonzero touch depth.
Never set on an FBA row --- no code path computes an FBA analogue.

\paragraph{Total book depth (C).} Mean of (bid depth $+$ ask
depth)$/2$ across the \emph{whole} book, every price level both sides
--- contrast with depth at best's touch-only figure.

\paragraph{Effective spread (U).}
\begin{equation}
  \mathrm{dev}_{\mathrm{bps}}(p_k, m_k, D_k) =
  \begin{cases}
    2(p_k - m_k)/m_k \cdot 10^{4} & D_k = \text{buy}\\
    -2(p_k - m_k)/m_k \cdot 10^{4} & D_k = \text{sell}\\
    2|p_k - m_k|/m_k \cdot 10^{4} & D_k \text{ unknown (FBA)}
  \end{cases}
\end{equation}
quantity-weighted per bucket, $m_k$ the pre-trade reference price (CDA:
book midpoint; FBA: the batch's own clearing price). \emph{None} if no
trade in the bucket had a usable reference price. The FBA has no
taker/maker distinction in a uniform-price batch, hence the third case.

\textbf{This third case is not a cosmetic detail.} Every FBA trade uses
the unsigned branch, so \emph{every} FBA contribution to this metric is
$\ge 0$ by construction, while a CDA contribution can be negative (a fill
better than the reference price). The two engines' means are therefore
not the same statistic: CDA's is a genuine signed average that can
partly cancel across trades; FBA's is a mean absolute deviation, which
cannot cancel and is mechanically at least as large as what a signed
average of the same underlying moves would show. Comparing the two
levels directly, as \Cref{ch:results} does, is still informative but
should not be read as isolating a pure mechanism effect on cost --- part
of any gap is this measurement asymmetry, not necessarily the trading
cost participants faced.

\paragraph{Realized spread (U), $\Delta \in \{1,5,30\}$\,s.} Same trades
and gate as effective spread, but against the reference price at
$t_k+\Delta$ instead of before the trade, found by a binary search over
the combined, time-sorted series of CDA midpoints and FBA clearing
prices (a book snapshot wins a tie against a batch at the identical
timestamp). \emph{None} for a horizon with no trade whose markout
falls within the retained window. The signed/unsigned asymmetry above
matters more here than for effective spread: because $D_k$ is
\emph{always} unknown for the FBA, its realized spread is a mean
\emph{absolute} price displacement over the next $\Delta$ seconds, not a
directional markout. For a price process that moves diffusively, the
expected absolute displacement grows roughly with $\sqrt{\Delta}$, so an
FBA realized spread that roughly doubles from $\Delta=1$s to $5$s and
roughly quintuples by $30$s (as \Cref{ch:results} reports) is consistent
with ordinary price diffusion over the markout window, not with a
mechanism-specific effect on adverse selection --- the CDA's signed
version, by contrast, can and does go negative, which is the
economically distinct claim that liquidity providers are on average
losing to informed flow after the fact. Reading the two as directly
comparable adverse-selection measures is the most likely place to
over-interpret this chapter's numbers.

\paragraph{Price impact (U), $\Delta \in \{1,5,30\}$\,s.} Effective
spread minus realized spread at the matching horizon --- the
adverse-selection component of the effective spread \emph{for the CDA,
where both terms are signed}; for the FBA this difference inherits the
mean-absolute-deviation character of both its inputs, so a negative FBA
value at a longer horizon means the mean absolute price move outgrew the
mean absolute submission-time deviation, not that adverse selection
reversed sign. Assigned only if
both terms are computable for that bucket and horizon.

\paragraph{Amihud illiquidity (U).} From the same combined reference-price
series, $\mathrm{close}$ is the last price observed in a bucket, carried
forward across empty buckets. For a bucket with a close and a computable
prior close $>0$:
\begin{equation}
  \mathrm{ret} = \frac{\mathrm{close}-\mathrm{prev\_close}}{\mathrm{prev\_close}},
  \qquad
  \mathrm{amihud} = \frac{|\mathrm{ret}|}{\mathrm{dollar\_volume}} \cdot 10^{6}
\end{equation}
if the bucket traded nonzero volume, else \emph{None} --- a real price
move on zero volume reports \emph{None}, not $\infty$.

\section{RQ2.2: Price Discovery Metrics}
\label{sec:formularq22}

\paragraph{Realized volatility (U).} From the reference-price series,
per bucket, consecutive-pair returns $r_i = (p_{i+1}-p_i)/p_i$
(a pair with a non-positive base price is skipped); the reported value
is $\sqrt{\sum_i r_i^2}$ --- the root sum of squares, the standard
realized-variance estimator, deliberately \emph{not} a standard
deviation about the mean. A single return in the bucket reports
$|r|$ directly. \emph{None} with fewer than two usable points.

\paragraph{Intra-interval price dispersion (U).}
From bucketed trade prices, $\mathrm{disp} =
\mathrm{stddev}(\{p_k\})/\mathrm{mean}(\{p_k\}) \cdot 10^{4}$ (population
standard deviation). This is the one column that departs from the
\emph{None}-versus-zero convention stated in \Cref{app:formulas}'s first
section: it reports $0.0$, not \emph{None}, whenever the bucket had at
least one trade but the mean price was non-positive. \emph{None} only
when the bucket had no trades at all. Approximately zero for the FBA by
construction (one clearing price per batch): empirically, exactly $0.00$
in the large majority but not literally all FBA-side buckets
(\Cref{ch:results} reports the exact share for the run behind this
thesis) --- a bucket occasionally contains trades from more than one
batch clear, most plausibly around a checkpoint-resumed run's seam
(\Cref{sec:experiments}), which this appendix states rather than
asserting a stronger guarantee than the implementation actually gives.

\paragraph{Kyle's lambda (U, but built differently per engine).} Per
bucket, the OLS slope through the origin, $\lambda = \sum_i x_i
y_i/\sum_i x_i^2$, in bps per unit of the base asset (not
$\mathrm{PRICE\_SCALE}$-denominated). CDA: one observation per
\emph{sweep} (same-timestamp trades sharing a side and a pre-trade
midpoint, grouped in timestamp order); $x$ is the net signed executed
quantity of the sweep, $y = (m_{\mathrm{post}}-m_{\mathrm{pre}})/
m_{\mathrm{pre}} \cdot 10^{4}$ with $m_{\mathrm{post}}$ the midpoint
5\,s after the sweep --- a fixed horizon, independent of the three
realized-spread horizons above. FBA: one observation per priced batch,
walked in strictly ascending order; $x$ is the batch's net order flow
(submitted buy quantity minus submitted sell quantity, captured before
price selection), $y$ the percentage change in $\clr$ from the previous
priced batch. \emph{None} if $\sum_i x_i^2 = 0$ for the bucket. On this
dataset the CDA value is positive-skewed (buy-side sweeps lift the
midpoint, on average); the FBA value clusters near zero --- a
consequence of the price-selection rule described in
\Cref{sec:environment}, which anchors to the previous clearing price and
is close to impervious to a single batch's flow imbalance, not a defect
in the estimator. \Cref{ch:results} reports both.

\paragraph{Pricing error (X, always empty in this dataset).}
$|\clr - p^{\mathrm{ref}}|$ or $|\midp - p^{\mathrm{ref}}|$ against an
external reference price, in bps. No code path in this codebase ever
assigns this column a value: the December~2025 SOL order-status archive
used here carries no oracle or mark-price feed, and none is
reconstructed elsewhere in the pipeline. Every row of both timeseries
CSVs has this field empty, for the entire month, on both engines --- a
documented, structural gap rather than a result of zero.

\section{RQ2.3 and Engine-Control Metrics}

\paragraph{Executed volume, executed notional (U).} Plain running sums
over the trade stream: quantity, and quantity $\times$ price in raw
$\mathrm{PRICE\_SCALE}$ units (summed before any division).

\paragraph{VWAP (U).} $\mathrm{executed\_notional}/\mathrm{executed\_volume}$
if the bucket executed nonzero volume, else \emph{None}.

\paragraph{Trade count (U).} A plain count, not \emph{None}-valued,
incremented once per trade.

\paragraph{Fill rate (U).} Bucketed by \emph{each order's own submission
bucket}, not the trade's bucket: $\sum \mathrm{filled\_qty} / \sum
\mathrm{orig\_qty}$ over orders first seen in that bucket. \emph{None} if
no order was first seen in the bucket. Because this thesis's interval is
1\,second and most orders take longer than that to reach their eventual
fill state, this bucketing choice --- not a low overall fill propensity
--- is the main driver of a low average value; \Cref{ch:results} reports
it alongside average time to execution, below, for that reason.

\paragraph{Average time to execution (U).}
For every order with a first fill, $(\text{first\_fill\_ts} -
\text{first\_seen\_ts})/10^{9}$ seconds, bucketed by first-seen time;
mean per bucket, \emph{None} if no order first seen in the bucket ever
received a fill.

\paragraph{Trader surplus (U, a plain value, not \emph{None}-valued).}
Accumulated per trade and per side independently: the buy side (when its
limit price is known) contributes $\max(0, \pi - p_k) \cdot q_k$, the
sell side $\max(0, p_k - \pi) \cdot q_k$ --- clipped at zero, so an
adverse fill never contributes a negative amount. This is realized
price improvement against each side's own submitted limit, in raw
$\mathrm{PRICE\_SCALE}$ units.

\paragraph{Order size inflation (U).} Bucketed by first-seen time,
aggregated per (bucket, participant) as (submitted total, filled total);
a participant contributes $\mathrm{submitted}/\mathrm{filled}$ to the
bucket's mean only when both are strictly positive --- an order that
never filled at all is excluded, deliberately, rather than folded in
through a floor on the denominator, which would let never-traded orders
dominate the statistic rather than genuinely oversized ones.
\emph{None} if no participant in the bucket had a computable ratio.

\paragraph{Order-to-trade ratio (U).} $\mathrm{bucket\_message\_count} /
\mathrm{trade\_count}$; the numerator counts \emph{every} message,
accepted and rejected, so its population does not line up one-to-one
with fill rate's. \emph{None} if the bucket executed no trades.

\paragraph{Boundary concentration (F).} Share of order arrivals landing
in the final 10\% of an FBA batch's own window; a zero-width window
contributes nothing. Numerator and denominator accumulate across every
batch closing in the bucket before dividing. \emph{None} if the bucket
had zero total messages across its batches.

\paragraph{Throughput, clearing latency (U, wall-clock and
non-deterministic across runs and machines).} Throughput:
$\mathrm{bucket\_message\_count} / \sum(\text{measured compute time})$.
Clearing latency: $\sum(\text{measured compute time})/\mathrm{count}$,
where the count is \emph{per submission} for the CDA (one book snapshot,
one timing, per actionable record) but \emph{per batch clear} for the
FBA (one timing per batch) --- the same column name measures two
different granularities of operation depending on engine, a caveat that
matters when \Cref{ch:results} compares the two directly. Neither column
is a claim about market quality; both exist to rule out ``one engine is
just slower'' as an alternative explanation for a measured difference
elsewhere.

\paragraph{Unexecuted residual share (F).} Summed across every batch
closing in the bucket: numerator $\sum |\,\mathrm{demand}(\clr) -
\mathrm{supply}(\clr)\,|$, denominator $\sum
\max(\mathrm{demand}(\clr),\mathrm{supply}(\clr))$. This is the
structural one-sidedness of the schedule \emph{at the clearing price}
within a batch, not the complement of fill rate and not ``the share of
submitted volume left unfilled'' in the everyday sense --- a batch can
match every unit it possibly can ($\min(\mathrm{demand},\mathrm{supply})$
executed in full) and still report a high value here if demand and
supply were far apart in size at $\clr$, which is common when $\tau$ is
short and a batch contains few orders. \emph{None} if the denominator is
zero.

\section{Known Gaps}

\begin{itemize}
  \item \textbf{Pricing error} is always empty for the reason stated
  above --- no oracle or mark-price feed in this dataset.
  \item \textbf{Throughput and clearing latency} are wall-clock and vary
  run to run and machine to machine; clearing latency additionally
  compares different operation granularities across engines (per
  submission for the CDA, per batch for the FBA).
  \item \textbf{Price impact} combines two separately quantity-weighted
  trade populations (effective spread's and realized spread's), which
  can differ for trades near the end of the retained window --- harmless
  in a full run, where every flushed bucket's markouts already exist by
  the time it is computed.
  \item \textbf{Kyle's lambda}'s differing CDA/FBA construction and its
  fixed 5\,s CDA horizon (independent of the three realized-spread
  horizons) are deliberate methodological choices, not a defect; the
  near-zero FBA result is itself a finding about the price-selection
  rule's insensitivity to flow imbalance, not noise.
  \item \textbf{Not computed in this codebase}: comparison of simulated
  CDA output against the trades Hyperliquid's own engine actually
  produced from the same flow, and the $\tau \to 0$ degenerate-parameter
  convergence check --- both require infrastructure or a parameter sweep
  this thesis did not build or run (\Cref{sec:validation}).
\end{itemize}
